% !TeX program = xelatex
\documentclass[aspectratio=169]{beamer}

\usepackage{import}

\subimport{../../}{system.tex}
\UseTemplateSet{
  layout = beamer,
  globalfonts = {cmu,shanggu},
  fonts = {hindi,sanskrit},
  features = {tables}
}

% Common packages (same as original preamble intent).
\usepackage{amsmath,amssymb}
\usepackage{tikz}

% Optional: ToC title (override per deck).
\providecommand{\AgendaTitle}{Contents}
\RenewDocumentCommand{\SA}{+m}{{\sanskritfont #1}}

\newcolumntype{V}{%
    @{}
        P{0.6cm} P{1cm} P{1.8cm}
        @{\hspace{1em}}
        P{0.6cm} P{1cm} P{1.8cm}
        @{\hspace{1em}}
        P{0.6cm} P{1cm} P{1.8cm}
    @{}
}

\title{梵語讀書會~第二講}
\author{\texorpdfstring{\SA{आकाश याङ्ग}}{आकाश याङ्ग}}
\date{March 13, 2026}


\begin{document}
\begin{frame}[plain]
  \titlepage
\end{frame}
\section{複習}

\begin{frame}{天城體拼讀練習}
    \begin{columns}
        \begin{column}{0.5\textwidth}
            \SA{तपःस्वाध्यायनिरतं तपस्वी वाग्विदां वरम् ।

                नारदं परिपप्रच्छ वाल्मीकिर्मुनिपुंगवम् ॥ १ ॥

                \vspace{1em}

                को न्वस्मिन्साम्प्रतं लोके गुणवान्कश्च वीर्यवान् ।

                धर्मज्ञश्च कृतज्ञश्च सत्यवाक्यो दृढव्रतः ॥ २ ॥

                \vspace{1em}

                चारित्रेण च को युक्तः सर्वभूतेषु को हितः ।

                विद्वान्कः कः समर्थश्च कश्चैकप्रियदर्शनः ॥ ३ ॥
            }
        \end{column}

        \begin{column}{0.5\textwidth}
            \SA{आत्मवान्को जितक्रोधो द्युतिमान्कोऽनसूयकः ।

                कस्य बिभ्यति देवाश्च जातरोषस्य संयुगे ॥ ४ ॥

                \vspace{1em}

                एतदिच्छाम्यहं श्रोतुं परं कौतूहलं हि मे ।

                महर्षे त्वं समर्थोऽसि ज्ञातुमेवंविधं नरम् ॥ ५ ॥

                \vspace{1em}

                श्रुत्वा चैतत्त्रिलोकज्ञो वाल्मीकेर्नारदो वचः ।

                श्रूयतामिति चामन्त्र्य प्रहृष्टो वाक्यमब्रवीत् ॥ ६ ॥
            }
        \end{column}
    \end{columns}
\end{frame}

\begin{frame}{複習}
    \begin{block}{\SA{अ}-語幹陽性名詞主賓格}
        \begin{NiceBooktable}{llll}{}{}
            \NiceHead{& \textbf{單數} & \textbf{雙數} & \textbf{複數} \\}
            \textbf{主格} & \SA{बालः} bālaḥ & \SA{बालौ} bālau & \SA{बालाः} bālāḥ \\
            \textbf{賓格} & \SA{बालम्} bālam & \SA{बालौ} bālau & \SA{बालान्} bālān \\
        \end{NiceBooktable}
    \end{block}

    \begin{block}{\SA{अ}-語幹中性名詞主賓格}
        \begin{NiceBooktable}{llll}{}{}
            \NiceHead{& \textbf{單數} & \textbf{雙數} & \textbf{複數} \\}
            \textbf{主格} & \SA{फलम्} phalam & \SA{फले} phale & \SA{फलानि} phalāni \\
            \textbf{賓格} & \SA{फलम्} phalam & \SA{फले} phale & \SA{फलानि} phalāni \\
        \end{NiceBooktable}
    \end{block}
\end{frame}

\begin{frame}{複習}
    \begin{block}{現在時\SA{परस्मैपद}第三人稱變位}
        \begin{NiceBooktable}{llll}{}{}
            \NiceHead{& \textbf{單數} & \textbf{雙數} & \textbf{複數} \\}
        \textbf{第三人稱} & \SA{गच्छति} gacchati & \SA{गच्छतः} gacchataḥ & \SA{गच्छन्ति} gacchanti \\
        \end{NiceBooktable}
    \end{block}

    \begin{block}{口語}
        \begin{itemize}
            \item \SA{नमो नमः। नमस्कारः। नमस्ते।}
            \item \SA{भवान्कः। भवती का। स कः। सा का। तत्किम्।}
            \item \SA{अहं रामः। अहं भारतदेशीयः। अहं सीता। अहं नेपलदेशीया।}
            \item \SA{स शिक्षकः। सा शिक्षका। तत्पुस्तकम्। }
        \end{itemize}
    \end{block}
\end{frame}

\begin{frame}{詞彙}
    \footnotesize
    \begin{block}{動詞}
        \begin{VocabCols}{3}
            \SA{आ-ह्वे} 1. \SA{आह्वयति} 叫過來
            
            \SA{उप-विश्} 6. \SA{उपविशति} 坐下
            
            \SA{क्रीड्} 1. \SA{क्रीडति} 玩耍
            
            \SA{क्रुध्} 4. \SA{क्रुध्यति} 發怒，生氣
            
            \SA{खाद्} 1. \SA{खादति} 喫
            
            \SA{गम्} 1. \SA{गच्छति} 走，去
            
            \SA{दृश्} 4. \SA{पश्यति} 看，看見
            
            \SA{नम्} 1. \SA{नमति} 問候，鞠躬，禮敬
            
            \SA{पठ्} 1. \SA{पठति} 誦，讀，學習
            
            \SA{पठ्} 10. \SA{पाठयति} 教
            
            \SA{पत्} 1. \SA{पतति} 落下，摔倒，飛
            
            \SA{प्रविश्} 6. \SA{प्रविशति} 走進，踏入
            
            \SA{श्रम्} 4. \SA{श्राम्यति} 疲倦
            
            \SA{सद्} 1. \SA{सीदति} 坐
            
            \SA{हस्} 1. \SA{हसति} 笑
        \end{VocabCols}
    \end{block}

    \begin{block}{名詞}
        \begin{VocabCols}{4}
            \SA{गोपाल} m. 瞿波羅
            
            \SA{फल} n. 水果，果實
            
            \SA{बाल} m. 孩童，男孩

            \SA{राम} m. 羅摩

            \SA{विद्यालय} m. 學校

            \SA{शिक्षक} m. 老師

            \SA{शिष्य} m. 學生

            \SA{सूक्त} n. 聖歌，讚詩
        \end{VocabCols}
    \end{block}

    \begin{block}{不變詞}
        \begin{VocabCols}{4}
            \SA{अत्र} 這裏，這兒

            \SA{अधुना} 現在

            \SA{ततः} 然後，從那裏

            \SA{तत्र} 那裏，那兒

            \SA{तदनु} 隨後，然後，接着

            \SA{पूर्वम्} 首先，以前

            \SA{सहसा} 突然，一下子，立刻
        \end{VocabCols}
    \end{block}
\end{frame}

\begin{frame}{翻譯參考}
    \begin{quotation}
        \SA{बालः शिक्षकः पश्यति। बालः शिष्यः। बालः शिक्षकं नमति। शिक्षक उपविशति। अधुना शिक्षकः सूक्तानि पठति। तदनु शुष्याः सूक्तानि पठन्ति। द्वौ शिष्यौ श्राम्यतः। शिक्षकः शिष्यौ पश्यति। शिक्षकः क्रुध्यति। तदनु शिष्याः क्रीडन्ति। सहसा बालः पतति। रामः पश्यति हसति। शिक्षकः क्रुध्यति॥}

        \vspace{1em}

        \rmfamily\itshape
        男孩看見老師。男孩是學生。男孩問候老師。老師坐下。現在老師誦讀聖歌。然後學生們誦讀聖歌。兩個學生疲倦了。老師看到這兩個學生。老師生氣了。然後學生們玩耍。突然一個男孩絆倒了。羅摩看見，笑了。老師生氣了。
    \end{quotation}
\end{frame}

\section{語法}
\begin{frame}{動詞現在時詞幹的構成}
    \small
    在梵語中，動詞構成現在時詞幹的方法有十類（\SA{गण}），其中第一、四、六、十類是構幹類（thematic），相對較簡單。
    \begin{block}{第一類動詞現在時詞幹的構成}
        \begin{itemize}
            \item 詞根元音使用\SA{गुण}，添加後綴\SA{अ}在詞根上\\
                \begin{tabular}{V}
                    \SA{यज्} & \SA{यजति} & 祭祀 
                    & \SA{रुह्} & \SA{रोहति} & 生長，登上
                    & \SA{स्मृ} & \SA{स्मरति} & 想起 \\
                \end{tabular}            
            \item 有尾輔音且去掉尾輔音後屬於長音節的詞根，元音不使用\SA{गुण}\\
                \begin{tabular}{V}
                    \SA{क्रीड्} & \SA{क्रीडति} & 玩耍 &
                    \SA{निन्द्} & \SA{निन्दति} & 斥責 \\
                \end{tabular}            
            \item 詞根若以\SA{इ}、\SA{ई}、\SA{उ}、\SA{ऊ}結尾，在變爲\SA{गुण}後需要進行詞內連聲（internal sandhi）\\
                \begin{tabular}{V}
                    \SA{जि} & \SA{जयति} & 勝利 &
                    \SA{नी} & \SA{नयति} & 帶領 &
                    \SA{भू} & \SA{भवति} & 是，成爲 \\
                \end{tabular}
            \item 有些第一類動詞現在時詞幹的構成並不規則
                \begin{tabular}{V}
                    \SA{गम्} & \SA{गच्छति} & 走 &
                    \SA{स्था} & \SA{तिष्ठति} & 站立 &
                    \SA{यम्} & \SA{यच्छति} & 給 \\
                    \SA{पा} & \SA{पिबति} & 喝 \\
                \end{tabular}
        \end{itemize}
    \end{block}
\end{frame}

\begin{frame}{動詞現在時詞幹的構成}
    \small
    \begin{block}{第四類動詞現在時詞幹的構成}
        \begin{itemize}

        \item 在詞根上直接加上後綴\SA{य}\\
            \begin{tabular}{V}
                \SA{स्निह्} & \SA{स्निह्यति} & 愛 &
                \SA{तुष्} & \SA{तुष्यति} & 滿意 &
                \SA{नृत्} & \SA{नृत्यति} & 跳舞 \\
            \end{tabular}
        \end{itemize}
    \end{block}

    \begin{block}{第六類動詞現在時詞幹的構成}
        \begin{itemize}
            \item 在詞根上直接加上後綴\SA{अ} \\
            \begin{tabular}{V}
                \SA{विश्} & \SA{विशति} & 進入 &
                \SA{तुद्} & \SA{तुदति} & 打，擊 &
                \SA{सृज्} & \SA{सृजति} & 創造 \\
            \end{tabular}
            \item 部分詞根存在不規則變化 \\
            \begin{tabular}{V}
                \SA{मृ} & \SA{म्रियते} & 死 &
                \SA{कॄ} & \SA{किरति} & 撒開，散佈 &
                \SA{सिच्} & \SA{सञ्चति} & 倒出，澆注 \\
            \end{tabular}
        \end{itemize}
    \end{block}
\end{frame}

\begin{frame}{動詞現在時詞幹的構成}
    \small
    \begin{block}{第十類動詞現在時詞幹的構成}
        \begin{itemize}
            \item 詞根有一個尾輔音，詞根元音是\SA{इ}或\SA{उ}，詞根元音使用\SA{गुण}，加上後綴\SA{अय}\\
            \begin{tabular}{V}
                \SA{घुष्} & \SA{घोषयति} & 宣告，叫嚷 \\
            \end{tabular}
            \item 詞根以簡單元音結尾，或詞根元音是\SA{अ}且有一個尾輔音，詞根元音使用\SA{वृद्धि}，加上後綴\SA{अय}\\
            \begin{tabular}{V}
                \SA{तड्} & \SA{ताडयति} & 打，擊 \\
            \end{tabular}
            \item 其他情況下，直接加上後綴\SA{अय} \\
            \begin{tabular}{V}
                \SA{चिन्त्} & \SA{चिन्तयति} & 考慮，想 &
                \SA{पूज्} & \SA{पूजयति} & 禮敬 \\
            \end{tabular}
        \end{itemize}
    \end{block}

    第十類動詞詞幹的構成可以應用到其他動詞詞幹上表示致使式，如：\\
    \hspace{2em}\begin{tabular}{V}
        \SA{पत्} & \SA{पातयति} & 砍倒 &
        \SA{पठ्} & \SA{पाठयति} & 教，講授 &
        \SA{दृश्} & \SA{दर्शयति} & 指示 \\
    \end{tabular}
\end{frame}

\begin{frame}{動詞現在時變位}
    \small
    動詞\SA{गम्} 「走，去」的現在時\SA{परस्मैपद}變位如下：

        \begin{NiceBooktable}{llll}{}{}
            \NiceHead{& \textbf{單數} & \textbf{雙數} & \textbf{複數} \\}
            \textbf{第一人稱} & \SA{गच्छामि} & \SA{गच्छावः} & \SA{गच्छामः} \\
            \textbf{第二人稱} & \SA{गच्छसि} & \SA{गच्छथः} & \SA{गच्छथ} \\ 
            \textbf{第三人稱} & \SA{गच्छति} & \SA{गच्छतः} & \SA{गच्छन्ति} \\
        \end{NiceBooktable}
    動詞\SA{लभ्} 「得到」的現在時\SA{आत्मनेपद}變位如下：

        \begin{NiceBooktable}{llll}{}{}
            \NiceHead{& \textbf{單數} & \textbf{雙數} & \textbf{複數} \\}
            \textbf{第一人稱} & \SA{लभे} & \SA{लभावहे} & \SA{लभामहे} \\
            \textbf{第二人稱} & \SA{लभसे} & \SA{लभेथे} & \SA{लभेध्वे} \\ 
            \textbf{第三人稱} & \SA{लभते} & \SA{लभेते} & \SA{लभन्ते} \\
        \end{NiceBooktable}
\end{frame}

\begin{frame}{\SA{अ}-語幹陽中性名詞變格}
    \begin{columns}
    \begin{column}{0.5\textwidth}
        以\SA{अ}結尾陽性名詞\SA{बाल}「男孩」變格如下：

        \begin{NiceBooktable}{llll}{}{}
            \NiceHead{& \textbf{單數} & \textbf{雙數} & \textbf{複數} \\}
            \textbf{主格} & \SA{बालः} & \SA{बालौ} & \SA{बालाः} \\
            \textbf{賓格} & \SA{बालम्} & \SA{बालौ} & \SA{बालान्} \\
            \textbf{具格} & \SA{बालेन} & \SA{बालाभ्याम्} & \SA{बालैः} \\
            \textbf{爲格} & \SA{बालाय} & \SA{बालाभ्याम्} & \SA{बालेभ्यः} \\
            \textbf{從格} & \SA{बालात्} & \SA{बालाभ्याम्} & \SA{बालेभ्यः} \\
            \textbf{屬格} & \SA{बालस्य} & \SA{बालयोः} & \SA{बालानाम्} \\
            \textbf{處格} & \SA{बाले} & \SA{बालयोः} & \SA{बालेषु} \\
            \textbf{呼格} & \SA{बाल} & \SA{बालौ} & \SA{बालाः} \\
        \end{NiceBooktable}
    \end{column}

    \begin{column}{0.5\textwidth}
        以\SA{अ}結尾中性名詞\SA{फल}「果子」變格如下：

        \begin{NiceBooktable}{llll}{}{}
            \NiceHead{& \textbf{單數} & \textbf{雙數} & \textbf{複數} \\}
            \textbf{主格} & \SA{फलम्} & \SA{फले} & \SA{फलानि} \\
            \textbf{賓格} & \SA{फलम्} & \SA{फले} & \SA{फलानि} \\
            \textbf{具格} & \SA{फलेन} & \SA{फलाभ्याम्} & \SA{फलैः} \\
            \textbf{爲格} & \SA{फलाय} & \SA{फलाभ्याम्} & \SA{फलेभ्यः} \\
            \textbf{從格} & \SA{फलात्} & \SA{फलाभ्याम्} & \SA{फलेभ्यः} \\
            \textbf{屬格} & \SA{फलस्य} & \SA{फलयोः} & \SA{फलानाम्} \\
            \textbf{處格} & \SA{फले} & \SA{फलयोः} & \SA{फलेषु} \\
            \textbf{呼格} & \SA{फलम्} & \SA{फले} & \SA{फलानि} \\
        \end{NiceBooktable}
    \end{column}
    \end{columns}
\end{frame}

\begin{frame}{名詞格的用法}
    \begin{block}{具格（Instrumental）}
        \begin{itemize}
            \item 表示行爲的工具或媒介\\
            \begin{quotation}\rmfamily
                \SA{रामः कुन्तेन मारयति।} 羅摩用矛殺。
            \end{quotation}
            \item 表示陪伴者，等同於加上後置詞\SA{सह}「和……一起」\\
            \begin{quotation}\rmfamily
                \SA{रामो बालेन (सह) गृहं गच्छति।} 羅摩和男孩一起回家。
            \end{quotation}
            \item 用於後置詞\SA{विना}「沒有」前\\
            \begin{quotation}\rmfamily
                \SA{रामो बालेन विना गृहं गच्छति।} 羅摩沒和男孩一起回家。
            \end{quotation}
            \item 在被動句中表示施動者
        \end{itemize}
    \end{block}
\end{frame}

\begin{frame}{名詞格的用法}
    \begin{block}{爲格（Dative）}
        \begin{itemize}
            \item 表示及物動詞的間接賓語\\
            \begin{quotation}
                \rmfamily
                \SA{रामो बालाय फलं यच्छति।} 羅摩給男孩一個水果。
            \end{quotation}
            \item 表達行爲的目的\\
            \begin{quotation}
                \rmfamily
                \SA{रामः पुत्राय देवं यजति।} 羅摩爲了兒子祭祀神。\\
                \SA{फलाय नगरं गच्छामि।} 我爲了水果去城裏。
            \end{quotation}
            \item 與表示「生氣」或「渴望」的動詞連用表示對象\\
            \begin{quotation}
                \rmfamily
                \SA{रामः पुत्राय क्रुध्यति।} 羅摩對兒子生氣。
            \end{quotation}
        \end{itemize}
    \end{block}
\end{frame}

\begin{frame}{名詞格的用法}
    \begin{block}{從格（Ablative）}
        \begin{itemize}
            \item 表示動作或移動過程的出發點\\
            \begin{quotation}
                \rmfamily
                \SA{रामो गृहादागच्छति।} 羅摩從家來。\\
                \SA{रामात्पुत्रः फलं लभते।} 兒子從羅摩那裏得到一個水果。
            \end{quotation}
            \item 表示行爲的原因\\
            \begin{quotation}
                \rmfamily
                \SA{क्रोधात्पुत्रं ताडयति।} 出於憤怒他打了兒子。
            \end{quotation}
            \item 與表示「害怕、保護、放棄」的動詞連用，表示來源\\
            \begin{quotation}
                \rmfamily
                \SA{देवो नरान्दुःखात्तारयति।} 神從苦難中救渡人們。
            \end{quotation}
            \item 一系列不變詞支配從格
        \end{itemize}
    \end{block}
\end{frame}

\begin{frame}{名詞格的用法}
    \begin{block}{屬格（Genitive）}
        \begin{itemize}
            \item 修飾名詞表示限定、所屬\\
            \begin{quotation}
                \rmfamily
                \SA{रामस्य पुत्रः पठति।} 羅摩的兒子學習。
            \end{quotation}
            \item 與動詞\SA{अस्}「是，存在」搭配表示「有」\\
            \begin{quotation}
                \rmfamily
                \SA{रामस्य पुत्रोऽस्ति।} 羅摩有一個兒子。
            \end{quotation}
            \item 一系列不變詞支配屬格
        \end{itemize}
    \end{block}
\end{frame}

\begin{frame}{名詞格的用法}
    \begin{block}{處格（Locative）}
        \begin{itemize}
            \item 表示行爲的地點\\
            \begin{quotation}
                \rmfamily
                \SA{रामो नगरे वसति।} 羅摩住在城裏。
            \end{quotation}
            \item 表示行爲的時間\\
            \begin{quotation}
                \rmfamily
                \SA{वसन्ते रामः क्षेत्रं गच्छति।} 羅摩在春天到田地去。
            \end{quotation}
            \item 與「投擲、擊中」等含義的動詞連用表示目標或方向\\
            \begin{quotation}
                \rmfamily
                \SA{रामः कुन्तं नरे क्षिपति।} 羅摩向那個人投擲矛。
            \end{quotation}
            \item 表示情感等的關聯性\\
            \begin{quotation}
                \rmfamily
                \SA{रामः पुत्रे स्निह्यति।} 羅摩愛兒子。
            \end{quotation}
            \item 用在形容詞最高級結構裏
        \end{itemize}
    \end{block}
\end{frame}

\begin{frame}{名詞格的用法}
    \begin{block}{呼格（Vocative）}
        \begin{itemize}
            \item 表示稱呼，常與\SA{}「喂」連用\\
            \begin{quotation}
                \rmfamily
                \SA{हे राम। कुत्र गच्छसि।} 喂，羅摩，你去哪兒？
            \end{quotation}
        \end{itemize}
    \end{block}

    名詞格的用法可以用下列例句總結：

    \begin{quotation}\rmfamily
        \SA{वसन्ते रामस्य पुत्रोऽश्वेन युद्धाय गृहात्क्षेत्रं गच्छति।}

        羅摩的兒子在春天裏爲了戰鬥從家出發騎着馬去戰場。
    \end{quotation}
\end{frame}

\begin{frame}{賓格的特殊用法}
    賓格還有以下特殊用法：
    \begin{itemize}
        \item 雙賓格：表示想和說的動詞，可以支配雙賓格\\
        \begin{quotation}
            \rmfamily
            \SA{रामो मार्गं बालं पृच्छति।} 羅摩向男孩問路。
        \end{quotation}
        \item 副詞賓格：形容詞的賓格單數可以相當於副詞\\
        \begin{quotation}
            \rmfamily
            \SA{नृपः सुखं जीवति।} 國王快樂地生活。
        \end{quotation}
        \item 時間賓格：時間名詞的賓格表示一段時間\\
        \begin{quotation}
            \rmfamily
            \SA{नृपः सहस्राणि वर्षाणि जीलति।} 國王活一千年。
        \end{quotation}
    \end{itemize}
\end{frame}

\begin{frame}{第三人稱陽中性代詞變格}
    \begin{columns}
        \begin{column}{0.5\textwidth}
            第三人稱陽性代詞\SA{सः}變格如下：

            \begin{NiceBooktable}{llll}{}{}
                \NiceHead{& \textbf{單數} & \textbf{雙數} & \textbf{複數} \\}
                \textbf{主格} & \SA{सः} & \SA{तौ} & \SA{ते} \\
                \textbf{賓格} & \SA{तम्} & \SA{तौ} & \SA{तान्} \\
                \textbf{具格} & \SA{तेन} & \SA{ताभ्याम्} & \SA{तैः} \\
                \textbf{爲格} & \SA{तस्मै} & \SA{ताभ्याम्} & \SA{तेभ्यः} \\
                \textbf{從格} & \SA{तस्मात्} & \SA{ताभ्याम्} & \SA{तेभ्यः} \\
                \textbf{屬格} & \SA{तस्य} & \SA{तयोः} & \SA{तेषाम्} \\
                \textbf{處格} & \SA{तस्मिन्} & \SA{तयोः} & \SA{तेषु} \\
            \end{NiceBooktable}
        \end{column}

        \begin{column}{0.5\textwidth}
            第三人稱中性代詞\SA{तद्}變格如下：

            \begin{NiceBooktable}{llll}{}{}
                \NiceHead{& \textbf{單數} & \textbf{雙數} & \textbf{複數} \\}
                \textbf{主格} & \SA{तत्} & \SA{ते} & \SA{तानि} \\
                \textbf{賓格} & \SA{तत्} & \SA{ते} & \SA{तानि} \\
                \textbf{具格} & \SA{तेन} & \SA{ताभ्याम्} & \SA{तैः} \\
                \textbf{爲格} & \SA{तस्मै} & \SA{ताभ्याम्} & \SA{तेभ्यः} \\
                \textbf{從格} & \SA{तस्मात्} & \SA{ताभ्याम्} & \SA{तेभ्यः} \\
                \textbf{屬格} & \SA{तस्य} & \SA{तयोः} & \SA{तेषाम्} \\
                \textbf{處格} & \SA{तस्मिन्} & \SA{तयोः} & \SA{तेषु} \\
            \end{NiceBooktable}
        \end{column}
    \end{columns}
\end{frame}

\begin{frame}{第三人稱代詞的作用}
    第三人稱代詞有以下作用：

        \begin{itemize}
            \item 用作人稱代詞\\
            \begin{quotation}
                \rmfamily
                \SA{स गृहं गच्छति।} 他回家。
            \end{quotation}
            \item 用作指示代詞\\
            \begin{quotation}
                \rmfamily
                \SA{स नरो गृहं गच्छति।} 那個人走進房子裏。
            \end{quotation}
            \item 用作定冠詞\\
            \begin{quotation}
                \rmfamily
                \SA{रामो बालं पश्यति। तं बालं ह्वयति। तदा स बाल आगच्छति।}\\
                羅摩看見男孩。他叫那個男孩。然後那個男孩過來。
            \end{quotation}
        \end{itemize}
\end{frame}

\begin{frame}{疑問代詞和不定代詞}
    \begin{block}{疑問代詞}
        陽性疑問代詞爲\SA{कः}，中性爲\SA{किम्}，與第三人稱代詞同樣變格。
        
        \begin{quotation}\rmfamily
            \SA{कः पठति।} 誰讀？\\
            \SA{बालः कं पठति।} 男孩看見誰？\\
            \SA{बालः किं खादति।} 男孩喫什麼？
        \end{quotation}
    \end{block}

    \begin{block}{不定代詞}
        不定代詞\SA{एक}意爲「一個，唯一一個，某一個」，形式上相當於形容詞，變格同第三人稱代詞。

        \begin{quotation}
            \rmfamily
            \SA{एकं पुस्तकं पठामि।} 我讀一本書。
        \end{quotation}
    \end{block}
\end{frame}

\begin{frame}{不變詞\SA{च}、\SA{न}、\SA{अपि}、\SA{इति}的用法}
    \begin{block}{\SA{च}的用法}
        \begin{itemize}
            \item 並列的詞通過\SA{च}連接，出現在每個詞後，或者只出現在最後\\
            \begin{quotation}
                \rmfamily
                \SA{बालौ रामं (च) गोपालं च ह्वयतः।} 兩個男孩呼喚羅摩和瞿波羅。
            \end{quotation}
            \item 連接兩個句子時，\SA{च}出現在第二個句子的第一個詞後面\\
            \begin{quotation}
                \rmfamily
                \SA{बालो गृहं गच्छति तत्र च रामं नमति।} 男孩回到家，並在那問候羅摩。
            \end{quotation}
        \end{itemize}
    \end{block}

    \begin{block}{\SA{न}的用法}
        \begin{itemize}
            \item \SA{न}放在動詞前或句首，表示句子的否定\\
            \begin{quotation}
                \rmfamily
                \SA{बालः फलं न खादति।} 男孩沒喫水果。
            \end{quotation}
        \end{itemize}
    \end{block}
\end{frame}

\begin{frame}{不變詞\SA{च}、\SA{न}、\SA{अपि}、\SA{इति}的用法}
    \begin{block}{\SA{अपि}的用法}
        \begin{itemize}
            \item \SA{अपि}放於所限定的詞後，用於表達「也，還，甚至，即使」\\
            \begin{quotation}
                \rmfamily
                \SA{रामोऽपि फलं खादति।} 羅摩也喫了一個水果。\\
                \SA{रामः फलमपि खादति।} 羅摩還喫了一個水果。\\
            \end{quotation}
            \item \SA{किम्}或\SA{अपि}置於陳述句首可以用來表達是非問句\\
            \begin{quotation}
                \rmfamily
                \SA{किं/अपि विद्यालयं गच्छसि।} 你去學校嗎？
            \end{quotation}
        \end{itemize}
    \end{block}

    \begin{block}{\SA{इति}的用法}
        \begin{itemize}
            \item \SA{इति}放在引語末，可以表示直接引語\\
            \begin{quotation}
                \rmfamily
                \SA{क्षीरं न पिबामीति बालो वदति।} 「我不喝牛奶」，男孩說。\\
                \SA{बालो वदति क्षीरं न पिबामीति।} 男孩說：「我不喝牛奶。」
            \end{quotation}
        \end{itemize}
    \end{block}
\end{frame}

\section{閱讀}
\begin{frame}{詞彙}
    \footnotesize
    \begin{block}{動詞}
        \begin{VocabCols}{3}
            \SA{आ-गम्} 1. \SA{आगच्छति} 來
            
            \SA{आ-नी} 1. \SA{आनयति} 帶來，拿來
            
            \SA{इष्} 6. \SA{इच्छति} 希望，想要
            
            \SA{क्षल्} 10. \SA{क्षालयति} 洗
            
            \SA{पच्} 1. \SA{पचति} 烹煮
            
            \SA{पा} 1. \SA{पिबति} 喝
            
            \SA{पूज्} 10. \SA{पूजयति} 禮敬，禮拜，供養
            
            \SA{यज्} 1. \SA{यजति} 祭祀
            
            \SA{यम्} 1. \SA{यच्छति} 給予，呈獻，遞給
            
            \SA{स्था} 1. \SA{तिष्ठति} 站立
            
            \SA{स्मृ} 1. \SA{स्मरति} 想起，回憶
        \end{VocabCols}
    \end{block}

    \begin{block}{名詞}
        \begin{VocabCols}{3}
            \SA{अन्न} n. 飯，食物
            
            \SA{क्षीर} n. 奶，乳
            
            \SA{गृह} n./m. 家，房屋
            
            \SA{जल} n. 水
            
            \SA{देव} m. 神
            
            \SA{पाद} m. 腳
            
            \SA{पुत्र} m. 兒子
            
            \SA{मुख} n. 臉，嘴
            
            \SA{सेवक} m. 僕人，侍從
        \end{VocabCols}
    \end{block}

    \begin{block}{不變詞}
        \begin{VocabCols}{2}
            \SA{अपि} 也，還，甚至，即使
            
            \SA{अभितः} （acc+）周圍，圍繞
            
            \SA{कुत्र} 哪裏
            
            \SA{सद्यः} 立刻
        \end{VocabCols}
    \end{block}
\end{frame}

\begin{frame}{基礎課文}
    \begin{block}{課文1}
        \SA{रामो गृहं गच्छति। तत्र तिष्ठति बालः। स रामं पश्यति। बालः सेवकानाह्वयति। सद्यः सेवका आगच्छन्ति। ते रामं नमन्ति। रामो जलमिच्छति। सेवको जलं यच्छति। पूर्वं रामः पादौ मुखं च क्षालयति। ततः स गृहं प्रविशति। अधुना रामो देवं स्मरति। सद्यः पुत्रौ राममागच्छतः। रामो देवं यजति। सेवका अन्नं पचन्ति। पूर्वं फलानि क्षीरं च सेवक आनयति। तदनु सोऽन्नमानयति। रामोऽन्नं खादति। ततः स फलानि खादति। तदनु स क्षीरं पिबति। पुत्रौ न खादतः। तौ क्रीडतः॥}
    \end{block}

    回答問題：
    \raggedcolumns
    \begin{multicols}{2}
        \begin{enumerate}
            \item \SA{रामः कुत्र गच्छति।}
            \item \SA{रामः किमिच्छति।}
            \item \SA{को रामं पश्यति।}
            \item \SA{कान्बाल आह्वयति।}
            \item \SA{रामः कं स्मरति।}
            \item \SA{के देवं पूजयन्ति।}
            \item \SA{सेवकाः किं पचन्ति।}
        \end{enumerate}
    \end{multicols}
\end{frame}

\begin{frame}{詞彙}
    \footnotesize
    \begin{block}{動詞}
        \begin{VocabCols}{3}
            \SA{अस्} 4. \SA{अस्यति} 扔，投擲，射
            
            \SA{क्षिप्} 6. \SA{क्षिपति} 扔，投擲
            
            \SA{घुष्} 10. \SA{घोषयति} 宣佈，宣告
            
            \SA{जि} 1. \SA{जयति} 勝利，戰勝，征服
            
            \SA{तुष्} 4. \SA{तुष्यति} 高興，滿意
            
            \SA{प्र-यम्} 1. \SA{प्रयच्छति} 給予，給
            
            \SA{प्र-शंस्} 1. \SA{प्रशंसति} 稱讚，讚美
            
            \SA{रक्ष्} 1. \SA{रक्षति} 保護，看護
            
            \SA{वद्} 1. \SA{वदति} 說
        \end{VocabCols}
    \end{block}

    \begin{block}{名詞}
        \begin{VocabCols}{4}
            \SA{कुन्त} m. 矛
            
            \SA{क्षेत्र} n. 田地，田野，國土
            
            \SA{जय} m. 勝利
            
            \SA{दान} n. 贈物，布施
            
            \SA{दूत} m. 信使
            
            \SA{नगर} n. 城
            
            \SA{नृप} m. 國王
            
            \SA{पुस्तक} n. 書
            
            \SA{योध} m. 戰士，士兵
            
            \SA{वीर} m. 英雄
            
            \SA{शर} m. 箭
        \end{VocabCols}
    \end{block}

    \begin{block}{不變詞}
        \begin{VocabCols}{4}
            \SA{अद्य} 今天，現在
            
            \SA{न} 不
            
            \SA{यदा...तदा} 當……時
            
            \SA{यदि...तर्हि} 如果……那麼
        \end{VocabCols}
    \end{block}
\end{frame}

\begin{frame}{基礎課文}
    \begin{block}{課文2}
        \SA{नृपो योधानाह्वयति। नृपो योधाश्च देवान्यजन्ति। अद्य जयाम इति नृपो वदति। नृपो योधाश्च क्षेत्रं गच्छन्ति। तत्र योधाः पूर्वं शरानस्यन्ति। ततो नृपः कुन्तान्क्षिपति। नृपो जयति। नपो वीर इति योधा वदन्ति नृपं च प्रशंसन्ति। यदि देवन्स्मरथ तर्हि देवा रक्षन्तीति नृपो वदति। तदनु नृपो दानानि प्रयच्छति। योधास्तुष्यन्ति। दूता नगरं गच्छन्ति जयं च घोषयन्ति॥}
    \end{block}

    回答問題：
    \raggedcolumns
    \begin{multicols}{2}
        \begin{enumerate}
            \item \SA{के देवान्यजन्ति।}
            \item \SA{कः कुन्तान्क्षिपति।}
            \item \SA{के नृपं प्रशंसन्ति।}
            \item \SA{योधाः किमस्यन्ति।}
            \item \SA{दूताः किं घोषयन्ति।}
            \item \SA{नृपः किं प्रयच्छति।}
        \end{enumerate}
    \end{multicols}
\end{frame}

\begin{frame}{詞彙}
    \footnotesize
    \begin{block}{動詞}
        \begin{VocabCols}{3}
            \SA{कृष्} 6. \SA{कृषति} 耕，犁
            
            \SA{कृष्} 1. \SA{कर्षति} 拉
            
            \SA{चिन्त्} 10. \SA{चिन्तयति} 想，考慮
            
            \SA{दृश्} 10. \SA{दर्शयति} 指給看
            
            \SA{नी} 1. \SA{नयति} 帶領
            
            \SA{नृत्} 4. \SA{नृत्यति} 跳舞
            
            \SA{प्रच्छ्} 6. \SA{पृच्छति} 問
            
            \SA{प्र-स्था} 1. \SA{प्रतिष्ठते} 出發
            
            \SA{लिप्} 6. \SA{लिम्पति} 塗，抹
            
            \SA{सिच्} 6. \SA{सिञ्चति} 倒出，澆注
            
            \SA{स्पृश्} 6. \SA{स्पृशति} 觸碰
            
            \SA{ह्वे} 1. \SA{ह्वयति} 叫，招呼，召喚
        \end{VocabCols}
    \end{block}

    \begin{block}{名詞}
        \begin{VocabCols}{4}
            \SA{अश्व} m. 馬
            
            \SA{कृषक} m. 農夫
            
            \SA{चित्र} n. 圖畫
            
            \SA{धन} n. 錢，財富
            
            \SA{प्रासाद} m. 宮殿
            
            \SA{मार्ग} m. 路，道路，街
            
            \SA{मित्र} n./m. 朋友
            
            \SA{युद्ध} n. 戰鬥
            
            \SA{रथ} m. （戰）車
            
            \SA{लाङ्गल} n. 犁
            
            \SA{वन} n. 森林
            
            \SA{विष} n. 毒藥
            
            \SA{शस्त्र} n. 武器
            
            \SA{सुख} n. 快樂，舒適
            
            \SA{सैन्य} n. 軍隊
            
            \SA{स्व} （前綴）自己的
            
            \SA{हस्त} m. 手
        \end{VocabCols}
    \end{block}

    \begin{block}{不變詞}
        \begin{VocabCols}{3}
            \SA{आदाय} （acc+）帶着
            
            \SA{किमर्थम्} 爲什麼
            
            \SA{सह} （ins+）一起
        \end{VocabCols}
    \end{block}
\end{frame}

\begin{frame}{基礎課文}
    \begin{block}{課文3}
        \SA{रामोऽश्वेन प्रतिष्ठते। गोपालो रामं पश्यति पृच्छति च। कुत्र गच्छसीति। प्रासादं गच्छामीति रामो वदति। किमर्थं प्रासादं गच्छसीति पुत्रः पृच्छति। नृपो धनेन सुखमिच्छति। ततः स युद्धाय गच्छति। अहं नृपेण सह युद्धाय गच्छामीति रामो वदति। अहमपि युद्धाय गच्छामीति गोपालो वदति। किं शस्त्रमादाय युद्धाय गच्छमीति पुत्रः पृच्छति। रामश्चिन्तयति। तदा स्वपुत्रं कुन्तं दर्शयति। गोपालः सुखेन नृत्यति। स हस्ताभ्यां कुन्तं स्पृशति। तदनु रामः पुत्राय शस्त्रं प्रयच्छति। अधुना रामः पुत्रेण सह प्रतिष्ठते। एको योधस्ताभ्यां मार्गं दर्शयति। रामः स्वपुत्रं वनं नयति। ततो रामो गोपालश्च क्षेत्रं पश्यतः। तत्र कृषका लाङ्गलैः कृषन्ति। तदनु तौ नगरं प्रविशतः प्रासादं च गच्छतः। तत्र तौ योधान्पश्यतः। योधा रथैरागच्छन्ति। ते शरान्विषेण लिम्पन्ति। तदनु नृपो ह्वयति। योधाः प्रतिष्ठन्ते। रामः पुत्रश्चापि सैन्येन सह प्रतिष्ठेते॥}
    \end{block}

    \raggedcolumns
    \begin{multicols}{2}
        \begin{enumerate}
            \item \SA{केन सह रामः प्रतिष्ठते।}
            \item \SA{किमर्थं रामो नृपं गच्छति।}
            \item \SA{नृपः केन सुखमिच्छति।}
            \item \SA{रामः कस्मै कुन्तं दर्शयति।}
            \item \SA{पुत्रः केन कुन्तं स्पृशति।}
            \item \SA{कृषकाः केन कृषन्ति।}
            \item \SA{केन योधाः शरान्लिम्पन्ति।}
        \end{enumerate}
    \end{multicols}
\end{frame}

\begin{frame}{詞彙}
    \begin{block}{動詞}
        \begin{VocabCols}{2}
            \SA{अव-रुह्} 1. \SA{अवरोहति} 下（車、馬等）
            
            \SA{आ-क्रम्} 1./4. \SA{आक्रामति/आक्राम्यति} 攻擊，接近，走近
            
            \SA{आ-रभ्} 1. \SA{आरभते} 開始
            
            \SA{आ-रुह्} 1. \SA{आरोहति} 登上，騎上
            
            \SA{चल्} 1. \SA{चलति} 動，移動，行進
            
            \SA{धाव्} 1. \SA{धावति} 跑
            
            \SA{भाष्} 1. \SA{भाषते} 說
            
            \SA{भू} 1. \SA{भवति} 是，成爲
            
            \SA{मुच्} 6. \SA{मुञ्चति} 射，釋放，解脫
            
            \SA{याच्} 1. \SA{याचते} 乞求，請求，索要
            
            \SA{युध्} 4. \SA{युध्यते} 戰鬥
            
            \SA{लभ्} 1. \SA{लभते} 獲得，得到
            
            \SA{लुभ्} 4. \SA{लुभ्यति} 貪求，渴望（+dat/loc）
        \end{VocabCols}
    \end{block}
\end{frame}

\begin{frame}{詞彙}
    \begin{block}{名詞}
        \begin{VocabCols}{3}
            \SA{क्रोध} m. 憤怒，生氣
            
            \SA{गज} m. 大象
            
            \SA{जन} m. （集合單數或複數）人，人們，族群，民族
            
            \SA{दर्शन} n. 看，外觀
            
            \SA{दुःख} n. 苦，苦痛，煩惱
            
            \SA{देश} m. 地方，國家
            
            \SA{नर} m. 人，男人
            
            \SA{पर} m. 敵人，陌生人
            
            \SA{पाप} n. 惡，罪惡
            
            \SA{पुरुष} m. 人，男人
            
            \SA{प्रसाद} m. 恩惠
            
            \SA{बल} n. 力量
            
            \SA{भय} n. 害怕，恐懼
            
            \SA{यज्ञ} m. 祭祀
            
            \SA{रत्न} n. 珍寶，寶石
            
            \SA{लोभ} m. 貪婪
            
            \SA{वृक्ष} m. 樹
        \end{VocabCols}
    \end{block}

    \begin{block}{不變詞}
        \begin{VocabCols}{2}
            \SA{कुतः} 從哪裏，爲什麼
            
            \SA{तस्मात्} 因此，所以
            
            \SA{पश्चात्} 此後，從後面
            
            \SA{पुनर्} 又
            
            \SA{प्रति} （acc+）朝，向，對於
            
            \SA{हि} 因爲，就，確實（不放在句首）
        \end{VocabCols}
    \end{block}
\end{frame}

\begin{frame}{基礎課文}
    \small
    
    \begin{block}{課文4}
        \SA{नृपो रथाद्योधानाह्वयति। सद्यो नृपस्य योधा आक्रामन्ति युध्यन्ते च। ते शरान्मुञ्चन्ति बलेन च कुन्तान्क्षिपन्ति। ततो नृपः स्वयोधाश्च जयन्ति। पश्चान्नृपो योधान्क्षेत्रान्नगरं नयति। सहसा योधानामश्वाः श्राम्यन्ति। तस्मादश्वा योधेभ्यो जलं लभन्ते। तदन्वश्वाः पुनश्चलन्ति। पूर्वं नृपस्य दूता नगरं प्रविशन्ति। ते नृपस्य जयं घोषयन्ति। नगरस्य जना गृहेभ्यो मार्गं प्रति धावन्ति। नृपस्य दर्शनाय जना आगच्छन्ति। अधुना नृपो योधाश्च नगरं प्रविशन्ति। तदनु जना नृपं प्रशंसन्ति। नृपो हि देशं परेभ्यो रक्षति। वृक्षेभ्यो नगरस्य बाला मार्गं पश्यन्ति। कुतो वृक्षानारोहथेत्येकः पुरुषो बालान्पृच्छति। गजस्य भयाद्वृक्षानारोहाम इति बाला वदन्ति। रामः पुत्रेण सह नगरं प्रविशति। ततो नृपोऽश्वादवरोहति प्रासादं च प्रविशति। पश्चात्स प्रासादाज्जनान्पश्यति। अधुना योधा दानाय लुभ्यन्ति। ते नृपाद्धनानि रत्नानि चेच्छन्ति। तस्मान्नृपं दानानि याचन्ते। ततो नृपो भाषते। लोभाच्च क्रोधाच्च दुःखं भवति। देवा नरान्दुःखाद्रक्षन्ति पापाच्च मुञ्चन्ति। देवानां प्रसादेन जयामः। तस्मात्पूर्वं देवान्यजामः। पश्चाद्दानानि यच्छामीति। ततो नृपो योधाश्च यज्ञमारभन्ते। जनास्तं यज्ञं पश्यन्ति। तदनु योधा नृपाद्दानानि लभन्ते॥}
    \end{block}

    \noindent
    \begin{minipage}[t]{0.27\textwidth}
        \begin{enumerate}
            \item \SA{कुतो नृपो योधानाह्वयति।}
            \item \SA{कस्याश्वाः श्राम्यन्ति।}
        \end{enumerate}
    \end{minipage}\hfill
    \begin{minipage}[t]{0.40\textwidth}
        \begin{enumerate}
            \setcounter{enumi}{2}
            \item \SA{दूताः किं घोषयन्ति।}
            \item \SA{नगरस्य जनाः कुतो मार्गं प्रति धावन्ति।}
        \end{enumerate}
    \end{minipage}\hfill
    \begin{minipage}[t]{0.27\textwidth}
        \begin{enumerate}
            \setcounter{enumi}{4}
            \item \SA{कस्मान्नृपो जनान्पश्यति।}
            \item \SA{कस्माद्दुःखं भवति।}
        \end{enumerate}
    \end{minipage}
\end{frame}

\begin{frame}{詞彙}
    \begin{block}{動詞}
        \begin{VocabCols}{2}
            \SA{आ-हृ} 1. \SA{आहरति} 采摘，收集
            
            \SA{ईक्ष्} 1. \SA{ईक्षते} 看，注視
            
            \SA{उद्-स्था} 1. \SA{उत्तिष्ठति} 站起，起來
            
            \SA{जप्} 1. \SA{जपति} （輕聲）念誦
            
            \SA{डी} 4./1. \SA{डीयते/डयते} 飛
            
            \SA{भ्रम्} 1./4. \SA{भ्रमति/भ्राम्यति} 漫遊
            
            \SA{रुह्} 1. \SA{रोहति} 生長
            
            \SA{वस्} 1. \SA{वसति} 居住，生活
            
            \SA{वि-कस्} 1. \SA{विकसति} 開花
            
            \SA{वृत्} 1. \SA{वर्तते} 位於，有，存在
        \end{VocabCols}
    \end{block}
\end{frame}

\begin{frame}{詞彙}
    \begin{block}{名詞}
        \begin{VocabCols}{3}
            \SA{आकाश} m./n. 天空
            
            \SA{आश्रम} m./n. 隱修處
            
            \SA{उद्यान} n. 花園
            
            \SA{कट} m. 墊子
            
            \SA{कन्दुक} m. 球
            
            \SA{कमल} n. 蓮花
            
            \SA{कुसुम} n. 花
            
            \SA{कैलास} m. 凱拉薩，山名
            
            \SA{ग्राम} m. 村莊
            
            \SA{तापस} m. 苦行者
            
            \SA{पर्वत} m. 山
            
            \SA{पुष्प} n. 花，花朵
            
            \SA{मध्य} m./n. 中間
            
            \SA{मन्त्र} m. 聖語，真言，咒語
            
            \SA{मृग} m. 野獸，鹿，羚羊
            
            \SA{वसन्त} m. 春天
            
            \SA{वानर} m. 猴子
            
            \SA{विहग} m. 鳥
            
            \SA{शिखर} m./n. 山峰，頂點
            
            \SA{शिव} m. 溼婆，神名
            
            \SA{ह्रद} m. 池塘
        \end{VocabCols}
    \end{block}

    \begin{block}{不變詞}
        \begin{VocabCols}{3}
            \SA{किंन्तु} 但是
            
            \SA{क्व} 哪裏
            
            \SA{प्रतिदिनम्} 每天
            
            \SA{प्रातर्} 早上
            
            \SA{विना} （ins+）沒有，不帶
            
            \SA{समीपम्} （gen+）在……附近
            
            \SA{समीपे} （gen+）在……附近
            
            \SA{हे} 喂，嗨，嘿
        \end{VocabCols}
    \end{block}
\end{frame}

\begin{frame}{基礎課文}
    \small
    
    \begin{block}{課文5}
        \SA{देशे नगराणि ग्रामाश्च वर्तन्ते। नगरेषु ग्रामेषु च जना वसन्ति। रामो नगरे न वसति किंतु ग्रामे वसति। रामस्य गृहं ग्रामस्य मध्ये वर्तते। प्रतिदिनं प्राता रामः शिवं पुष्पैः पूजयति। किंत्वद्य गृहे पुष्पाणि न वर्तन्ते। हे सेवक। पुष्पाणि कुत्र वर्तन्त इति रामः पृच्छति। सेवकः कुसुमेभ्य उद्यानं गच्छति। तत्र कुसुमानि वृक्षेषु विकसन्ति। उद्याने ह्रदोऽपि वर्तते। ह्रदे कमलानि भवन्ति। जलेन विना कमलानि न रोहन्ति। सेवको वृक्षात्पुष्पाण्याहरति। स कुसुमानि रामाय यच्छति। अधुना रामः कट उपविशति। रामः कुसुमैः शिवं पूजयति। स मन्त्राञ्जपति। पश्चाद्राम उत्तिष्ठति गृहाच्चोद्यानं गच्छति। गृहस्य समीपे पर्वता वर्तन्ते। रामः पर्वतान्पश्यति चिन्तयति च। पर्वतानां शिखरेषु देवा वसन्ति। कैलासस्य शिखरे शिवो वसति। यदा प्रतिदिनं शिवं पूजयामि तदा शिवस्तुष्यतीति। तदा स आकाशमीक्षते। आकाशे विहगा डयन्ते। रामस्य गृहस्य समीपे वनं वर्तते। वने मृगा वसन्ति। वने गजा भ्रमन्ति। तत्र रामो वृक्षेषु वानरानपीक्षते। रामस्य ग्राम एक आश्रमो वर्तते। आश्रमे तापसा वसन्ति। तेऽपि प्रातः शिवं नमन्ति॥}
    \end{block}

    \raggedcolumns
    \begin{multicols}{2}
        \begin{enumerate}
            \item \SA{रामः क्व वसति।}
            \item \SA{रामस्य गृहं कुत्र वर्तते।}
            \item \SA{कंं देवं रामः पूजयति।}
            \item \SA{देवाः कुत्र वसन्ति।}
            \item \SA{कस्य शिखरे शेवो वसति।}
            \item \SA{कमलानि कुत्र विकसन्ति।}
            \item \SA{मृगाः कुत्र वसन्ति।}
        \end{enumerate}
    \end{multicols}
\end{frame}

\section{寫作}
\begin{frame}{翻譯練習}
    \begin{block}{翻譯練習1}
        \rmfamily\itshape
        學生們回家。那裏僕人們做飯。兩個男孩進房間。他倆想要水。僕人立刻拿來水。然後兩個男孩洗腳。現在羅摩來了。羅摩叫僕人們過來。僕人們來了並且站着。羅摩現在想要食物。然後僕人們拿來食物。羅摩和兩個兒子喫。
    \end{block}

    \begin{block}{翻譯練習2}
        \rmfamily\itshape
        國王叫一個戰士過來。他問他：「馬兒們站在哪裏？」戰士說：「我把馬兒們來帶。」然後戰士們騎上馬。國王說：「我們現在去城裏。」然後國王和戰士們進城。首先國王投擲一支矛，然後戰士們射箭。國王征服了這座城。他給予賞賜，戰士們很高興。
    \end{block}
\end{frame}

\begin{frame}{翻譯練習}
    \begin{block}{翻譯練習3}
        \rmfamily\itshape
        今天瞿波羅沒有步行上學。僕人帶來馬。羅摩問：「你倆什麼時候去學校？」僕人說：「我倆現在去」然後瞿波羅和僕人騎馬去學校。瞿波羅以雙手向老師行禮並進了學校。沒有書學生們不學習。老師給學生們拿來書。然後學生們讀。現在老師給學生們看圖畫。學生們看圖畫。然後學生們回家。
    \end{block}

    \begin{block}{翻譯練習4}
        \rmfamily\itshape
        國王進入城市。在宮殿附近，他從馬上下來。人們從家裏出來，跑向宮殿。戰士們站在那裏，並讚美國王。突然來了一個信使說：「國王現在講話。」然後國王連同一個僕人來了，說：「今天我連同軍隊從這座城市出發。首先我們祭祀衆神。然後我們登上馬。如果我們勝利，我給戰士們賞賜。」戰士們立刻從房子裏取來武器。國王從馬上觀看戰士們的武器。然後國王帶着軍隊出城。人們從房屋中看國王的軍隊。
    \end{block}
\end{frame}

\begin{frame}{翻譯練習}
    \begin{block}{翻譯練習5}
        \rmfamily\itshape
        我是瞿波羅的朋友。我每天去瞿波羅家。在他家附近有一座花園。在花園中間有一條路。在那條路上我們玩球。花園裏也有很多樹。樹上坐着猴子們。他們驅走鳥兒們。然後他們喫樹的果實。我們從花園也看見山。瞿波羅說，在山的頂峰住着衆神。因此他每天禮敬衆神。
    \end{block}
\end{frame}

\end{document}
